\documentclass[10pt,fleqn]{article}
\usepackage{hyperref}
\usepackage{graphicx}


\setlength{\topmargin}{-.75in}
\addtolength{\textheight}{2.00in}
\setlength{\oddsidemargin}{.00in}
\addtolength{\textwidth}{.75in}

\nofiles

\pagestyle{empty}

\setlength{\parindent}{0in}

\input{/cygdrive/m/tex/commands/commands}

\begin{document}
%%%%%%%%%%%%%%%%%%%%%%%%%%%%%%%%%%%%%%%%%%%%%%%%%%%%%%%%%%%%%%%%%%%%%%%%%%%%%%%%
%%%%%%%%%%%%%%%%%%%%%%%%%%%%%%%%%%%%%%%%%%%%%%%%%%%%%%%%%%%%%%%%%%%%%%%%%%%%%%%%
\vskip0.1in\hrule\vskip0.1in \noindent
{\bf Math 4610 Fundamentals of Computational Mathematics  - Topic 17.}
\vskip0.1in\hrule\vskip0.1in \noindent
In this section we will focus on the finite limitations of number representation
on computers. Computer capabilities have advanced rapidly over the past decades.
Computer hardware and software use a finite number of bits or memory addresses
to represent numbers.

Over the years computer archetecture has increased this resolution by using more
bits
to account for smaller and larger numbers. Starting with 8-bit, 16-bit, and then
onto 32-bit and 64-bit numbers, the resolution has gotten better, but will never
be able to represent an infinite number of digits needed to store irrational
numbers.

Even a number like \(1/3\) must be truncated after a finite number of digits due
to the way computers store numbers using bits. Data is stored in a binary format
on all computers. So, the number \(1/3\) becomes
\[
  ({1\over 3}_{10} = (0.10101010101010\ldots)_2
\]
where the subscript refers to the base of the number. The representation
repeats, but still requires an infinite number of digits to exactly represent
the number.
%%%%%%%%%%%%%%%%%%%%%%%%%%%%%%%%%%%%%%%%%%%%%%%%%%%%%%%%%%%%%%%%%%%%%%%%%%%%%%%%
%%%%%%%%%%%%%%%%%%%%%%%%%%%%%%%%%%%%%%%%%%%%%%%%%%%%%%%%%%%%%%%%%%%%%%%%%%%%%%%%
\vskip0.1in\hrule\vskip0.1in \noindent
{\bf Approximation Using Floating Point Numbers.} 
\vskip0.1in\hrule\vskip0.1in \noindent
It would be nice to be able to determine how well an arbitrary real number can
be approximated before we do any other computations. In fact, we can determine
the accuracy coding a simple loop. The idea is to add successively smaller
numbers to a specific number that can be represented exactly in a computer
register. For example, the number, \(x=0\), is a number that can be represented
exactly in memory. Two other examples are \(x=1\) and \(x=2\). Powers of two up
to a point can be represented exactly. We will use \(x=1\). A loop written in
Java that will achieve our object is:
\vskip0.1in\hrule\vskip0.1in \noindent
\begin{verbatim}

      //
      // set up some variables
      // ---------------------
      //
      float one = 1.0f;
      float seps = 1.0f;
      float appone = 1.0f;
      //
      // loop over more than enough iterations to computer the precision
      // ---------------------------------------------------------------
      //
      for(int k=0; k<200) {
        //
        // compute the approximation
        // -------------------------
        //
        appone = one + seps;
        //
        // test the approximation
        // ----------------------
        //
        if(Math.abs(appone-one) == 0.0f) {
          printf("single precision machine epsilon = " + seps);
          return;
        }
        //
        // reduce the perturbation values
        // ------------------------------
        //
        seps = 0.5f * seps;
        //
      }

\end{verbatim}
\vskip0.1in\hrule\vskip0.1in \noindent
This code will print a message as soon as the difference between the exact value
of one and the approximation. There are a couple of notes that one should 
consider in the code. Note that the lines that initialize values to numerical
values need an additional \lq\lq f\rq\rq\ on the end of the numerical value. For
example,
\vskip0.1in\hrule\vskip0.1in \noindent
\begin{verbatim}

      float one = 1.0f;

\end{verbatim}
\vskip0.1in\hrule\vskip0.1in \noindent
This ensures that the numbers are cast to float, which might be considered as
single precision.

The same code in double prevision can be written using the following.
\vskip0.1in\hrule\vskip0.1in \noindent
\begin{verbatim}

      //
      // set up some variables
      // ---------------------
      //
      double one = 1.0;
      double seps = 1.0;
      double appone = 1.0;
      //
      //
      // loop over more than enough iterations to computer the precision
      // ---------------------------------------------------------------
      //
      for(int k=0; k<200) {
        //
        // compute the approximation
        // -------------------------
        //
        appone = one + seps;
        //
        // test the approximation
        // ----------------------
        //
        if(Math.abs(appone-one) == 0.0) {
          printf("single precision machine epsilon = " + seps);
          return;
        }
        //
        // reduce the perturbation values
        // ------------------------------
        //
        seps = 0.5 * seps;
        //
      }

\end{verbatim}
\vskip0.1in\hrule\vskip0.1in \noindent
Note the omission of the \lq\lq\ f\rq\rq\ at the end of the constant values.
%%%%%%%%%%%%%%%%%%%%%%%%%%%%%%%%%%%%%%%%%%%%%%%%%%%%%%%%%%%%%%%%%%%%%%%%%%%%%%%%
%%%%%%%%%%%%%%%%%%%%%%%%%%%%%%%%%%%%%%%%%%%%%%%%%%%%%%%%%%%%%%%%%%%%%%%%%%%%%%%%
\vskip0.1in\hrule\vskip0.1in \noindent
{\bf Floating Point Notation} 
\vskip0.1in\hrule\vskip0.1in \noindent
It is helpful to have some notation to define the exact value of a number and an
approximation. We will define the approximation of \(x\) as
\[
  x \apprx fl(x)
\]
We will measure the error via the error formula
\[
  \makebox{error } =  {{ | \apprx fl(x) - x | }\over{ | x | }}
\]
This gives a scaled value of the error. The scaling is determined by the size or
magnitude of the exact number being approximated. So, for example, if \(x=1/3\)
then the error will look like
\[
  \makebox{error } =  {{ | \apprx fl(1/3) - 1/3 | }\over{ | 1/3 | }}
\]
The floating point value of \(x=1/3\) is dependent on the computer and the
precision implemented in hardware and software on your computer.
%%%%%%%%%%%%%%%%%%%%%%%%%%%%%%%%%%%%%%%%%%%%%%%%%%%%%%%%%%%%%%%%%%%%%%%%%%%%%%%%
%%%%%%%%%%%%%%%%%%%%%%%%%%%%%%%%%%%%%%%%%%%%%%%%%%%%%%%%%%%%%%%%%%%%%%%%%%%%%%%%
\vskip0.1in\hrule\vskip0.1in \noindent
{\bf IEEE Standards for Single and Double Prevision} 
\vskip0.1in\hrule\vskip0.1in \noindent
To standardize preision across platforms and all devices that might do some
computational arithmetic, the Institute of Electrical and Electronics Engineers
(IEEE) have established standards that computers must adhere to so that results
from computing can be trusted across different computers produced by various
companies - for example, Dell, HP, and a myriad of other computer vendors. The
standards include both number representation and standards related to arithmetic
operations. If we define
\[
  \epsilon = {1\over 2}\ 2^{-t}
\]
%%%%%%%%%%%%%%%%%%%%%%%%%%%%%%%%%%%%%%%%%%%%%%%%%%%%%%%%%%%%%%%%%%%%%%%%%%%%%%%%
%%%%%%%%%%%%%%%%%%%%%%%%%%%%%%%%%%%%%%%%%%%%%%%%%%%%%%%%%%%%%%%%%%%%%%%%%%%%%%%%
\vskip0.1in\hrule\vskip0.1in \noindent
{\bf Absolute and Relative Error: Introduction.} 
\vskip0.1in\hrule\vskip0.1in \noindent
In most cases, the solution of a mathematical problem can only be approximated.
So, it would be a good idea to have a way to measure the error between the exact
solution and the approximation of that solution. We will define two types of
error. These are absolute error and relative error. The {\bf absolute error} is
the absolute value of the difference between the approximation and the exact
value for the solution. That is, if $x^*$ is the exact value approximated by
$x$, then
$$ e_{abs} = | x - x^{*} | $$
defines the absolute error. The {\bf relative error} is a scaled error defined
by
$$ e_{rel} = {{| x - x^{*} |}\over{|x^*|}} $$
So, the relative error is a scaled or percent error based on the magnitude of
the exact value.
%%%%%%%%%%%%%%%%%%%%%%%%%%%%%%%%%%%%%%%%%%%%%%%%%%%%%%%%%%%%%%%%%%%%%%%%%%%%%%%%
%%%%%%%%%%%%%%%%%%%%%%%%%%%%%%%%%%%%%%%%%%%%%%%%%%%%%%%%%%%%%%%%%%%%%%%%%%%%%%%%

%%%%%%%%%%%%%%%%%%%%%%%%%%%%%%%%%%%%%%%%%%%%%%%%%%%%%%%%%%%%%%%%%%%%%%%%%%%%%%%%
%%%%%%%%%%%%%%%%%%%%%%%%%%%%%%%%%%%%%%%%%%%%%%%%%%%%%%%%%%%%%%%%%%%%%%%%%%%%%%%%
\vskip0.1in\hrule\vskip0.1in \noindent
{\bf Absolute and Relative Error: Examples.} 
\vskip0.1in\hrule\vskip0.1in \noindent
If we are trying to find the roots of the polynomial
$$
  p(x) = x^5 + x^3 - 2\ x^2 + 5\ x
$$
we can see that $x=0$ is one solution. To find other roots we can use Newton's
method to generate a sequence of approximations given a starting point. That is,
we will generate a sequence
$$
   S = \left\lbrace x_k \right\rbrace_{k=0}^{\infty}
$$
Newton's method will be covered a bit later in this course. However, what we
will want is for the sequence to converge to a root, say $x^*$. This can be
rephrased as
$$ | x_k - x^* | \rightarrow 0 $$
which implies the absolute error will tend to zero as $k$ tends to $\infty$.
Using the relative error we want
$$ {{ | x_k - x^* | }\over{ | x^* | }} \rightarrow 0 $$
For the polynomial define in this example, there will be problems in using the
relative error near the zero roots. So, if the sequence starts to converg to the
zero root, we would need to use the absolute error as a measure.
%%%%%%%%%%%%%%%%%%%%%%%%%%%%%%%%%%%%%%%%%%%%%%%%%%%%%%%%%%%%%%%%%%%%%%%%%%%%%%%%
%%%%%%%%%%%%%%%%%%%%%%%%%%%%%%%%%%%%%%%%%%%%%%%%%%%%%%%%%%%%%%%%%%%%%%%%%%%%%%%%

%%%%%%%%%%%%%%%%%%%%%%%%%%%%%%%%%%%%%%%%%%%%%%%%%%%%%%%%%%%%%%%%%%%%%%%%%%%%%%%%
%%%%%%%%%%%%%%%%%%%%%%%%%%%%%%%%%%%%%%%%%%%%%%%%%%%%%%%%%%%%%%%%%%%%%%%%%%%%%%%%
\vskip0.1in\hrule\vskip0.1in \noindent
{\bf Absolute and Relative Error: A Numerical Example.} 
\vskip0.1in\hrule\vskip0.1in \noindent
As an illustration of how absolute and relative errors compare, we can consider
some numerical examples. The following table gives a pair of numbers along with
both the absolute and relative errors.

\begin{table}[h!]
   \caption{Absolute and Relative Error Values}
   \begin{center}
   \begin{tabular}{c|c|c|c}
     \hline
     \textbf{$x$} & \textbf{$x^*$} & \textbf{abs. err.} & \textbf{rel. err.} \\
     \hline
     0.01 & 0.1 & 0.09 & 0.9 \\
     \hline
     1.01 & 1.0 & 0.01 & 0.01 \\
     \hline
     2.0   & 3.0  & 1.0 & 0.5 \\
     \hline
     10.0  & 9.0  & 1.0 & 0.1 \\
     \hline
     100.0 & 99.0 & 1.0 & 0.01 \\
     \hline
   \end{tabular}
   \end{center}
\end{table}

Note that when both the approximation and exact value are close to zero, the
absolute error becomes a better measure of error than the relative error and
for large values (at least greater than one) one should expect that the better
measure of the error is the relative error. We will use both in the development
of algorithms to numerically solve mathematical problems.






The following is a list of sources for error that need to be taken into account
by compoutational scientists.
\begin{enumerate}
  \item {\bf Modeling Errors} These errors can occur when assumptions are made
        about the phenomena being studied. For example, one may consider a model
        of the solar system where the planets are assumed to be spheres, which
        is not the case.
  \item {\bf Measurement Errors:} These errors occur when instruments are used
        to measure physical quantities. For example, the temperature of molten
        lava might be measured to within one or two degrees based on the
        magnitude of the exact temperature. The fractional part of the 
        measurement would characterize the error.
  \item {\bf Discretization Error:} In order to computer solutions to
        mathematical problems using computers, it necessary that the model be
        finite and discrete. For example, weather models based on systems of
        partial differential equations require a discretization of the continuous model to fit in the
        discrete framework of a computer simulation.
\end{enumerate}
Note that Github will not allow you the luxury of creating empty folders. This
is an advantage in using \lq\lq git\rq\rq\ on a local machine. When changes are
\lq\lq pushed\rq\rq\ to Github, empty folders are ignored. So, let's get started
on the formatting the homework solutions portion of the repository for the
class.
%%%%%%%%%%%%%%%%%%%%%%%%%%%%%%%%%%%%%%%%%%%%%%%%%%%%%%%%%%%%%%%%%%%%%%%%%%%%%%%%
%%%%%%%%%%%%%%%%%%%%%%%%%%%%%%%%%%%%%%%%%%%%%%%%%%%%%%%%%%%%%%%%%%%%%%%%%%%%%%%%
\end{document}
